% Options for packages loaded elsewhere
\PassOptionsToPackage{unicode}{hyperref}
\PassOptionsToPackage{hyphens}{url}
\documentclass[
  ignorenonframetext,
]{beamer}
\newif\ifbibliography
\usepackage{pgfpages}
\setbeamertemplate{caption}[numbered]
\setbeamertemplate{caption label separator}{: }
\setbeamercolor{caption name}{fg=normal text.fg}
\beamertemplatenavigationsymbolsempty
% remove section numbering
\setbeamertemplate{part page}{
  \centering
  \begin{beamercolorbox}[sep=16pt,center]{part title}
    \usebeamerfont{part title}\insertpart\par
  \end{beamercolorbox}
}
\setbeamertemplate{section page}{
  \centering
  \begin{beamercolorbox}[sep=12pt,center]{section title}
    \usebeamerfont{section title}\insertsection\par
  \end{beamercolorbox}
}
\setbeamertemplate{subsection page}{
  \centering
  \begin{beamercolorbox}[sep=8pt,center]{subsection title}
    \usebeamerfont{subsection title}\insertsubsection\par
  \end{beamercolorbox}
}
% Prevent slide breaks in the middle of a paragraph
\widowpenalties 1 10000
\raggedbottom
\AtBeginPart{
  \frame{\partpage}
}
\AtBeginSection{
  \ifbibliography
  \else
    \frame{\sectionpage}
  \fi
}
\AtBeginSubsection{
  \frame{\subsectionpage}
}
\usepackage{iftex}
\ifPDFTeX
  \usepackage[T1]{fontenc}
  \usepackage[utf8]{inputenc}
  \usepackage{textcomp} % provide euro and other symbols
\else % if luatex or xetex
  \usepackage{unicode-math} % this also loads fontspec
  \defaultfontfeatures{Scale=MatchLowercase}
  \defaultfontfeatures[\rmfamily]{Ligatures=TeX,Scale=1}
\fi
\usepackage{lmodern}
\ifPDFTeX\else
  % xetex/luatex font selection
\fi
% Use upquote if available, for straight quotes in verbatim environments
\IfFileExists{upquote.sty}{\usepackage{upquote}}{}
\IfFileExists{microtype.sty}{% use microtype if available
  \usepackage[]{microtype}
  \UseMicrotypeSet[protrusion]{basicmath} % disable protrusion for tt fonts
}{}
\makeatletter
\@ifundefined{KOMAClassName}{% if non-KOMA class
  \IfFileExists{parskip.sty}{%
    \usepackage{parskip}
  }{% else
    \setlength{\parindent}{0pt}
    \setlength{\parskip}{6pt plus 2pt minus 1pt}}
}{% if KOMA class
  \KOMAoptions{parskip=half}}
\makeatother
\usepackage{longtable,booktabs,array}
\newcounter{none} % for unnumbered tables
\usepackage{calc} % for calculating minipage widths
\usepackage{caption}
% Make caption package work with longtable
\makeatletter
\def\fnum@table{\tablename~\thetable}
\makeatother
\setlength{\emergencystretch}{3em} % prevent overfull lines
\providecommand{\tightlist}{%
  \setlength{\itemsep}{0pt}\setlength{\parskip}{0pt}}
% Slide layout defaults for warning-clean scientific decks.
\usepackage{etoolbox}
\IfFileExists{xurl.sty}{\usepackage{xurl}}{}
\IfFileExists{seqsplit.sty}{\usepackage{seqsplit}}{\newcommand{\seqsplit}[1]{#1}}
\protected\def\breakseq#1{\seqsplit{#1}}
\protected\def\breaktt#1{\begingroup\ttfamily\seqsplit{#1}\endgroup}
\setlength{\emergencystretch}{6em}
\tolerance=5000
\hbadness=10000
\hfuzz=1pt
\setlength{\tabcolsep}{2pt}
\AtBeginEnvironment{longtable}{\tiny\renewcommand{\arraystretch}{0.86}\setlength{\tabcolsep}{1pt}}
\AtBeginEnvironment{tabular}{\tiny\renewcommand{\arraystretch}{0.86}\setlength{\tabcolsep}{1pt}}

% Natbib and cross-reference fallbacks — slides don't load natbib
% or cleveref, but combined-PDF manuscript prose may emit these
% commands. The fallback renders citations as a bracketed key list
% and cross-references as detokenized labels so slides don't fail on
% undefined control sequences or raw underscores. \providecommand is
% a no-op if packages load later.
\providecommand{\citep}[1]{[#1]}
\providecommand{\citet}[1]{#1}
\providecommand{\citealp}[1]{#1}
\providecommand{\citeauthor}[1]{#1}
\providecommand{\citeyear}[1]{#1}
\providecommand{\cref}[1]{\texttt{\detokenize{#1}}}
\providecommand{\Cref}[1]{\texttt{\detokenize{#1}}}
\usepackage{bookmark}
\IfFileExists{xurl.sty}{\usepackage{xurl}}{} % add URL line breaks if available
\urlstyle{same}
% fallback for those not using the hyperref driver hyperxmp:
\makeatletter
\@ifundefined{xmpquote}{\newcommand{\xmpquote}[1]{#1}}{}
\makeatother
\hypersetup{
  hidelinks,
  pdfcreator={LaTeX via pandoc}}

\author{\texorpdfstring{}{}}
\date{}

\begin{document}

\section{Security and Provenance}\label{security-and-provenance}

\begin{frame}[fragile,allowframebreaks]{Security and Provenance}
Research integrity requires more than reproducibility; it requires
verifiable authorship. In an era of generative AI, automated scraping,
and synthetic media, the ability to prove that a document was produced
by a specific pipeline at a specific time is a critical defense against
fabrication and misattribution. The W3C PROV data model
\breakseq{[@moreau2013provdm]} establishes a formal vocabulary for expressing
provenance records---entities, activities, and agents connected by
derivation, generation, and attribution relations. Digital watermarking,
pioneered by Cox et al. \breakseq{[@cox1997secure]} for multimedia integrity
verification, provides the foundational signal-processing theory for
embedding imperceptible provenance markers within artifacts.
\texttt{template/} implements a domain-specific provenance layer that
embeds these relations directly into the PDF artifact itself, using four
complementary steganographic and cryptographic mechanisms.
\end{frame}

\begin{frame}[allowframebreaks]{Threat Model}
\protect\phantomsection\label{threat-model}
The steganography subsystem defends against three classes of threats:

\begin{enumerate}
\tightlist
\item
  \textbf{Unauthorized redistribution}: A manuscript is scraped and
  republished without attribution. The steganographic watermark survives
  the redistribution and can be used to prove original authorship.
\item
  \textbf{Content tampering}: Figures or text are modified after
  publication. The SHA-256 hash embedded in the PDF metadata detects any
  modification to the file contents.
\item
  \textbf{Provenance forgery}: An attacker claims to have produced a
  document they did not author. The build timestamp, commit hash, and
  pipeline metadata embedded in the watermark provide a verifiable chain
  of custody.
\end{enumerate}
\end{frame}

\begin{frame}[fragile,allowframebreaks]{Steganographic Layers}
\protect\phantomsection\label{steganographic-layers}
The system applies four complementary layers of provenance information:

\begin{block}{Layer 1: PDF Metadata Injection}
\protect\phantomsection\label{layer-1-pdf-metadata-injection}
The \breaktt{inject\_pdf\_metadata} function writes structured metadata
into both the PDF Info dictionary and an XMP (Extensible Metadata
Platform) packet:

\begin{itemize}
\tightlist
\item
  \texttt{/Creator}: Pipeline identifier
\item
  \texttt{/Producer}: Module path
  (\breaktt{infrastructure.steganography})
\item
  \texttt{/CreationDate}: UTC timestamp in ISO 8601 format
\item
  \texttt{/Author}: From \texttt{config.yaml}
\item
  \texttt{/Title}: From \texttt{config.yaml}
\item
  Custom fields: DOI, ORCID, repository URL
\end{itemize}
\end{block}

\begin{block}{Layer 2: Cryptographic Hashing}
\protect\phantomsection\label{layer-2-cryptographic-hashing}
Before watermarking, a SHA-256 hash of the rendered PDF is computed and
stored in:

\begin{itemize}
\tightlist
\item
  The output manifest (\breaktt{output/manifest.json})
\item
  The PDF metadata (\texttt{/Subject} field)
\item
  An external hash file
  (\texttt{output/\textless{}name\textgreater{}.sha256})
\end{itemize}

This enables post-hoc verification: anyone with the hash can verify that
the PDF has not been modified since rendering.
\end{block}

\begin{block}{Layer 3: Alpha-Channel Text Overlay}
\protect\phantomsection\label{layer-3-alpha-channel-text-overlay}
A semi-transparent text overlay is applied to each page of the PDF,
encoding:

\begin{itemize}
\tightlist
\item
  Build timestamp
\item
  Git commit hash (short SHA)
\item
  Project name
\item
  Pipeline version
\end{itemize}

The overlay is rendered at low opacity (typically 3--5\% alpha) to be
invisible during normal viewing but detectable through image analysis.
It survives printing (as a faint watermark) and standard PDF operations.

A representative overlay text string takes the following form:

\begin{verbatim}
template/ | built: 2026-03-19T14:23:11Z | commit: a4f2c1b | pipeline: v2.0.0 | project: template
\end{verbatim}

This single line, tiled across each page at 3--5\% opacity, encodes the
complete build provenance chain: the system identifier, ISO 8601 build
timestamp, short Git commit hash, pipeline version, and project name.
Together these fields allow a verifier to reconstruct---from the
watermark alone---which version of the code, at which moment in time,
produced the document.
\end{block}

\begin{block}{Layer 4: QR Code Injection}
\protect\phantomsection\label{layer-4-qr-code-injection}
An optional QR code is generated containing a URL pointing to the
repository (e.g., \breaktt{github.com/docxology/template}). The QR code
is placed in a configurable position (default: bottom-right corner of
the last page) at a specified size.
\end{block}
\end{frame}

\begin{frame}[fragile,allowframebreaks]{The \texttt{secure\_run.sh}
Orchestrator}
\protect\phantomsection\label{the-secure_run.sh-orchestrator}
The steganographic pipeline is orchestrated by \texttt{secure\_run.sh},
a Bash script that wraps the standard \texttt{run.sh} pipeline with
post-processing steganography:

\begin{enumerate}
\tightlist
\item
  Execute the standard YAML-declared pipeline (14 stages; default full
  10) pipeline for the target project.
\item
  The \texttt{secure\_run.sh} script invokes
  \breaktt{SteganographyProcessor}.
\item
  Apply metadata injection, hashing, text overlay, and QR code
  injection.
\item
  Save the secured PDF alongside the original.
\item
  Generate a steganography report in JSON format.
\end{enumerate}

The orchestrator processes either a single specified project or all
discovered projects sequentially.
\end{frame}

\begin{frame}[allowframebreaks]{Tamper Detection}
\protect\phantomsection\label{tamper-detection}
Verification is performed by comparing the stored SHA-256 hash against a
freshly computed hash of the distributed PDF. Any modification---even a
single bit flip---produces a hash mismatch. The alpha-channel overlay
provides a secondary, visual verification channel that does not require
access to the original hash.
\end{frame}

\begin{frame}[fragile,allowframebreaks]{Limitations}
\protect\phantomsection\label{limitations}
\begin{itemize}
\tightlist
\item
  Alpha-channel overlays are stripped by some PDF optimization tools
  (e.g., \breaktt{qpdf\ -\/-optimize}).
\item
  QR codes are visible and may be removed by a determined attacker.
\item
  The system does not provide non-repudiation in the cryptographic
  sense---it does not use digital signatures with private keys. Future
  versions may integrate GPG or X.509 signing.
\item
  The provenance model is pipeline-centric rather than fully
  PROV-compliant. The path from \texttt{template/}'s current
  metadata-based provenance to full W3C PROV-compliant traces involves
  four steps: (1) \emph{entity identification}---assigning stable
  identifiers (URIs or SWHIDs) to each pipeline input (manuscript files,
  data, config); (2) \emph{activity logging}---recording each pipeline
  stage as a PROV Activity with start/end timestamps (already encoded in
  the watermark overlay); (3) \emph{agent attribution}---binding each
  Activity to the pipeline version and Git commit hash (already encoded
  in the overlay and PDF metadata); (4) \emph{PROV-O
  serialization}---emitting the provenance graph as OWL-RDF (PROV-O) or
  text (PROV-N) alongside the PDF. Steps (2) and (3) are already
  implemented; steps (1) and (4) are the primary remaining gaps. A
  future \breaktt{infrastructure.provenance} module would close both gaps
  automatically.
\end{itemize}
\end{frame}

\begin{frame}[fragile,allowframebreaks]{Relationship to Software Supply
Chain Integrity}
\protect\phantomsection\label{relationship-to-software-supply-chain-integrity}
The steganographic provenance layer operates at the \emph{document
level}---it certifies the integrity of a specific PDF artifact. A
complementary concern is \emph{build-level} provenance: certifying that
the pipeline itself was executed with verified source code and
dependencies. Frameworks such as in-toto \breakseq{[@torresarias2019intoto]}
and SLSA (Supply-chain Levels for Software Artifacts) address this
concern by defining attestation chains from source commit through build
steps to final artifact. The NTIA's minimum elements for a Software Bill
of Materials (SBOM) \breakseq{[@ntia2021sbom]} further standardize the
enumeration of software components and dependencies---essential for
establishing the provenance lineage of build environments. SLSA defines
four graduated levels of build integrity:

{\def\LTcaptype{none} % do not increment counter
\begin{longtable}[]{@{}
  >{\centering\arraybackslash}p{(\linewidth - 4\tabcolsep) * \real{0.2727}}
  >{\raggedright\arraybackslash}p{(\linewidth - 4\tabcolsep) * \real{0.2955}}
  >{\raggedright\arraybackslash}p{(\linewidth - 4\tabcolsep) * \real{0.4318}}@{}}
\toprule\noalign{}
\begin{minipage}[b]{\linewidth}\centering
SLSA Level
\end{minipage} & \begin{minipage}[b]{\linewidth}\raggedright
Requirement
\end{minipage} & \begin{minipage}[b]{\linewidth}\raggedright
\texttt{template/} Status
\end{minipage} \\
\midrule\noalign{}
\endhead
1 & Provenance document exists & ✓ SHA-256 manifest + steganographic
metadata \\
2 & Version-controlled build scripts & ✓ All scripts in git \\
3 & Isolated build environment & \textasciitilde{} Docker support exists
but not enforced in CI \\
4 & Hermetic, reproducible builds & N -- future work \\
\bottomrule\noalign{}
\end{longtable}
}

Future versions of \texttt{template/} may generate SLSA-compatible
provenance attestations alongside the steganographic watermarks,
creating a two-layer provenance model: in-toto attests that the build
pipeline was executed with the claimed source code, while the
steganographic layer attests that the PDF was produced by that pipeline
at a specific time.
\end{frame}

\begin{frame}[fragile,allowframebreaks]{Relationship to FAIR and Formal
Provenance Standards}
\protect\phantomsection\label{relationship-to-fair-and-formal-provenance-standards}
The steganographic layer supports the FAIR for Research Software
(FAIR4RS) principles \breakseq{[@barker2022fair4rs]} at the artifact level.
PDFs carry embedded metadata (Findability) in standardized XMP format
(Interoperability). The SHA-256 hash manifest enables persistent
integrity verification (a prerequisite for Reusability). The
Documentation Duality standard ensures that the software producing the
artifact is inspectable and well-documented (satisfying FAIRsoft
\breakseq{[@garijo2024fairsoft]} metadata and documentation indicators). Full
PROV-compliant provenance traces---capturing the derivation chain from
source data through analysis scripts to rendered PDF---are a natural
extension and a primary target for future development.

Software Heritage \breakseq{[@cosmo2020softwareheritage]} complements this
picture at the source-code level: by archiving the \texttt{template/}
repository and assigning a reproducible SWHID (Software Hash Identifier)
to each commit, Software Heritage makes the pipeline itself---not just
its output---a citable, persistent digital artifact. A published SWHID
alongside the PDF DOI creates a complete, two-artifact citation record:
the paper's content is versioned via DOI; the code that generated it is
versioned via SWHID. This combination satisfies the Katz et al.
\breakseq{[@katz2021software]} software citation principles' requirement that
software used in research be independently citable and permanently
accessible.
\end{frame}

\end{document}
